%% 9장 반올림 오차 해석
\chapter{반올림 오차 해석}
\label{ch:rounding-error}

\abstract{연산 하나의 반올림 오차는 단위 반올림 $u$로 묶이지만, 수많은 연산이 이어지는 실제 계산에서 그 오차들이 어떻게 쌓이고 증폭되는지는 따로 따져야 한다. 이 장은 그 방법론을 다룬다. 먼저 근삿값이 참값에서 얼마나 떨어졌는지를 재는 전진 오차와, 근삿값이 어떤 가까운 입력에 대한 정확한 답인지를 묻는 후진 오차를 구별하고, 문제 자체가 입력의 작은 변화를 얼마나 키우는지를 재는 조건수를 정의한다. 이 셋을 잇는 ``전진 오차 $\lesssim$ 조건수 $\times$ 후진 오차''가 이 장의 중심 원리이다. 이어서 반올림 인수의 곱을 하나로 묶는 $\gamma_n$ 기법으로 합과 내적의 오차를 증명하고, 합하는 순서를 나무 모양으로 바꾸면 오차의 상한이 $n$에서 $\log n$으로 줄어드는 이유를 본다. 마지막으로 소거가 정확도를 무너뜨리는 원리와 그것을 피하도록 식을 고쳐 쓰는 방법들을, 분산과 이차방정식과 순위 검정의 통계량을 예로 살핀다.}

\section{오차의 척도}
\label{sec:error-measures}

참값 $y$를 근삿값 $\hat y$로 계산했을 때 $|\hat y-y|$를 \term{절대오차}{absolute error}, $y\ne0$이면 $|\hat y-y|/|y|$를 \term{상대오차}{relative error}라 한다. 부동소수점 수의 반올림은 크기에 비례하는 오차를 만들므로(8장) 수치 계산에서 자연스러운 척도는 상대오차이며, 상대오차가 $10^{-k}$ 정도이면 대략 유효숫자 $k$자리가 맞는다. binary64에서 $u\approx1.1\times10^{-16}$이므로 한 번 반올림한 값은 유효숫자 약 $16$자리가 맞고, 상대오차가 $10^{-10}$인 결과는 그 가운데 $6$자리를 잃은 셈이다.

벡터나 행렬에서는 노름으로 재는 상대오차 $\|\hat y-y\|/\|y\|$와 성분마다 재는 $\max_i|\hat y_i-y_i|/|y_i|$가 크게 다를 수 있다. 노름으로 재면 큰 성분의 오차만 보이고 작은 성분은 완전히 틀려도 드러나지 않기 때문이다. 어느 척도로 정확도를 약속하는지는 문제의 성격에 따라 정해야 하며, 이 장에서는 주로 스칼라 결과를 다루되 필요하면 성분별 척도를 쓴다.

\section{전진 오차, 후진 오차, 조건수}
\label{sec:forward-backward}

\subsection{후진 오차}

함수 $f$의 값 $y=f(x)$를 어떤 알고리즘이 $\hat y$로 계산했다고 하자. $|\hat y-y|$를 재는 것은 가장 직접적인 질문이지만, 참값 $y$를 모르는 상황에서는 이 양을 직접 다루기 어렵다. 다른 질문을 던질 수도 있다. $\hat y$는 \textbf{어떤 입력에 대한} 정확한 답인가? $\hat y=f(x+\Delta x)$를 만족하는 $\Delta x$ 가운데 가장 작은 것의 상대 크기
\[
\eta(\hat y)=\min\bigl\{\,|\Delta x|/|x|:\ \hat y=f(x+\Delta x)\,\bigr\}
\]
를 $\hat y$의 \term{후진 오차}{backward error}라 하고, 이에 맞서 $|\hat y-y|/|y|$를 \term{전진 오차}{forward error}라 한다. 후진 오차가 $u$ 정도인 알고리즘을 \term{후진 안정}{backward stable}이라 한다. 그런 알고리즘이 낸 답은 입력을 반올림하는 것만큼의 변화를 준 입력에 대한 정확한 답이고, 입력 자체가 측정값이거나 이미 반올림된 값이라면 그보다 나은 답을 기대할 근거가 없다.

8장의 표준 모형이 가장 단순한 예이다. $a\oplus b=(a+b)(1+\delta)=a(1+\delta)+b(1+\delta)$이므로, 계산한 합은 각 입력을 상대적으로 $u$ 이하만큼 바꾼 입력의 정확한 합이다. 곧 부동소수점 덧셈은 후진 안정이다. 결과 자체의 상대오차, 곧 전진 오차가 얼마인지는 이 사실만으로 정해지지 않는데 다음의 조건수에 달려 있다.

\subsection{조건수}

입력의 작은 상대 변화가 출력을 상대적으로 얼마나 바꾸는지를 문제의 \term{조건수}{condition number}라 한다. 미분 가능한 스칼라 함수에서는
\[
\frac{f(x+\Delta x)-f(x)}{f(x)}\approx\frac{xf'(x)}{f(x)}\cdot\frac{\Delta x}x
\]
이므로 상대 조건수를
\begin{equation}
\label{eq:cond-scalar}
\kappa_f(x)=\left|\frac{xf'(x)}{f(x)}\right|
\end{equation}
로 정의한다. 조건수는 알고리즘이 아니라 \textbf{문제}의 성질이며, 어떤 알고리즘을 쓰든 입력에 있는 상대오차 $\epsilon$은 출력에서 $\kappa_f\epsilon$ 정도로 나타난다. 몇 가지 예를 보자. $f(x)=\sqrt x$는 $\kappa=1/2$로 입력의 오차를 오히려 줄인다. $f(x)=e^x$는 $\kappa=|x|$이므로 $x=700$ 근처에서는 입력의 마지막 비트 하나가 결과의 $700$배 오차가 된다. $f(x)=\log x$는 $\kappa=1/|\log x|$이어서 $x\approx1$에서 한없이 커진다. 그리고 두 수의 뺄셈 $f(a,b)=a-b$는 두 입력을 각각 상대적으로 $\epsilon$ 이하만큼 바꿀 때
\begin{equation}
\label{eq:cond-sub}
\frac{|(a+\Delta a)-(b+\Delta b)-(a-b)|}{|a-b|}\le\frac{|a|+|b|}{|a-b|}\,\epsilon
\end{equation}
이므로 조건수가 $(|a|+|b|)/|a-b|$이고, 가까운 두 수의 뺄셈에서 아주 크다.

\subsection{세 양의 관계}

후진 오차가 $\eta$인 근삿값은 $\hat y=f(x+\Delta x)$, $|\Delta x|\le\eta|x|$로 쓸 수 있으므로, 조건수의 정의에서 일차 근사로
\begin{equation}
\label{eq:rule-of-thumb}
\text{전진 오차}\ \lesssim\ \text{조건수}\times\text{후진 오차}
\end{equation}
이다. 이 부등식은 수치 해석 전체를 꿰는 원리이다. 정확도를 잃는 원인을 두 가지로 나누어, 하나는 알고리즘이 책임지는 후진 오차로, 다른 하나는 문제가 정하는 조건수로 돌리기 때문이다. 후진 안정인 알고리즘으로 조건수가 작은 문제를 풀면 결과가 정확하고, 조건수가 $10^{k}$인 문제를 후진 안정인 알고리즘으로 풀면 유효숫자 약 $k$자리를 잃는다. 이 손실은 알고리즘의 결함이 아니라 문제의 본성이므로, 더 정밀한 산술을 쓰거나 문제를 다시 정식화하지 않는 한 피할 수 없다.

4장에서 본 선형 연립방정식의 조건수 $\kappa(A)=\|A\|\|A^{-1}\|$이 이 원리의 대표적인 예이다. 부분 피벗 Gauss 소거는 실용적인 의미에서 후진 안정이므로(19장), 계산한 해의 상대오차는 대략 $\kappa(A)u$이다. 이 관계에서 행렬의 조건수가 $10^{12}$이면 해의 유효숫자는 $4$자리 정도밖에 믿을 수 없다는 경험칙이 나온다.

\begin{instardust}[\code{docs/design/htest.md} §3.1]
가설검정의 p값도 조건수의 언어로 정확도를 약속한다. 통계량 $s$에서 p값 $p(s)$를 계산할 때, 통계량 자체에 있는 오차는 $\kappa_p=|s\,p'(s)/p(s)|$배로 커져 p값에 나타난다. 분포의 깊은 꼬리에서는 이 값이 크다. 정규분포의 위쪽 꼬리 $p(s)\approx\varphi(s)/s$에서는 $\kappa_p\approx s^2$이므로, $s=10$이면 통계량의 상대오차가 p값에서는 $100$배가 된다. 그래서 stardust의 정확도 계약은 p값의 오차를 $u$가 아니라 $\kappa_p$를 곱한 단위로 적는다. 알고리즘이 책임질 수 있는 것은 후진 오차뿐이라는 \cref{eq:rule-of-thumb}의 생각을 계약의 형식으로 옮긴 것이다.
\end{instardust}

\section{반올림 오차의 누적}
\label{sec:accumulation}

\subsection{$\gamma_n$ 기법}

연산이 이어지면 반올림 인수 $(1+\delta_i)$들이 곱해지거나 나누어진다. 이것을 일일이 다루면 식이 금세 다루기 어려워지므로, 그 곱을 하나의 인수로 묶는 보조정리를 먼저 증명해 둔다\cite{higham2002accuracy}.

\begin{lemma}
\label{lem:gamma-n}
$|\delta_i|\le u$, $\rho_i\in\{1,-1\}$($i=1,\dots,n$)이고 $nu<1$이면
\[
\prod_{i=1}^n(1+\delta_i)^{\rho_i}=1+\theta_n,\qquad|\theta_n|\le\gamma_n:=\frac{nu}{1-nu}
\]
이다.
\end{lemma}

\begin{proof}
$n$에 대한 귀납법을 쓴다. $n=1$이면 $\rho_1=1$일 때 $|\theta_1|=|\delta_1|\le u\le\gamma_1$이고, $\rho_1=-1$일 때 $|\theta_1|=|\delta_1|/|1+\delta_1|\le u/(1-u)=\gamma_1$이다. $n-1$까지 성립한다고 하고 $1+\theta_n=(1+\theta_{n-1})(1+\delta_n)^{\rho_n}$으로 쓰자. $\rho_n=1$이면 $\theta_n=\theta_{n-1}+\delta_n+\theta_{n-1}\delta_n$이므로
\begin{align*}
|\theta_n|&\le\gamma_{n-1}+u+\gamma_{n-1}u=\frac{(n-1)u+u\bigl(1-(n-1)u\bigr)+(n-1)u^2}{1-(n-1)u}\\
&=\frac{nu}{1-(n-1)u}\le\frac{nu}{1-nu}
\end{align*}
이다. $\rho_n=-1$이면 $\theta_n=(\theta_{n-1}-\delta_n)/(1+\delta_n)$이므로
\[
|\theta_n|\le\frac{\gamma_{n-1}+u}{1-u}=\frac{(n-1)u+u(1-(n-1)u)}{(1-(n-1)u)(1-u)}\le\frac{nu}{1-nu}
\]
인데, 마지막 부등식은 분자가 $nu$ 이하이고 분모 $(1-(n-1)u)(1-u)=1-nu+(n-1)u^2$이 $1-nu$ 이상이기 때문이다.
\end{proof}

$nu\ll1$이면 $\gamma_n\approx nu$이므로, 반올림 $n$번의 효과는 많아야 대략 $nu$의 상대오차라고 읽으면 된다. 이 묶음은 다음 규칙들로 쉽게 이어 붙일 수 있는데 모두 정의에서 곧바로 확인된다.
\[
(1+\theta_m)(1+\theta_n)=1+\theta_{m+n},\qquad\gamma_m+\gamma_n+\gamma_m\gamma_n\le\gamma_{m+n},\qquad c\,\gamma_n\le\gamma_{cn}\ (c\ge1).
\]
binary64에서 $\gamma_{10^6}\approx1.11\times10^{-10}$이므로 백만 번의 반올림이 한 방향으로 쌓이는 최악의 경우에도 유효숫자 $10$자리는 남는다.

\subsection{합의 오차}

$s=x_1+\cdots+x_n$을 왼쪽부터 차례로 더하는 \term{재귀 합}{recursive summation} $\hat s_1=x_1$, $\hat s_i=\hat s_{i-1}\oplus x_i$를 생각하자. 표준 모형을 펼치면
\begin{align*}
\hat s_n&=x_1(1+\delta_2)\cdots(1+\delta_n)+x_2(1+\delta_2)\cdots(1+\delta_n)\\
&\quad+x_3(1+\delta_3)\cdots(1+\delta_n)+\cdots+x_n(1+\delta_n)
\end{align*}
이고, 항 $x_i$에 붙은 인수는 많아야 $n-1$개이다. \cref{lem:gamma-n}을 쓰면 다음을 얻는다.

\begin{theorem}
\label{thm:recursive-sum}
재귀 합은
\[
\hat s_n=\sum_{i=1}^nx_i(1+\theta^{(i)}),\qquad|\theta^{(i)}|\le\gamma_{n-1}
\]
을 만족한다. 따라서 재귀 합은 후진 안정이고, 전진 오차는
\[
|\hat s_n-s|\le\gamma_{n-1}\sum_{i=1}^n|x_i|
\]
로 묶인다.
\end{theorem}

전진 오차를 $|s|$로 나누면 상대오차가 $\gamma_{n-1}\cdot\sum|x_i|/|\sum x_i|$ 이하인데, 뒤의 몫이 바로 합의 조건수이다. 모든 항의 부호가 같으면 조건수는 $1$이고 상대오차는 $\gamma_{n-1}$ 이하이지만, 항들이 서로 지워져 합이 작아지면 조건수가 커지고 상대오차의 상한도 커진다. 이 상한은 최악의 경우를 가정한 것이라는 점도 기억해 두자. 반올림 오차들이 서로 독립인 무작위 양처럼 행동한다면 그 합은 $\sqrt n$에 비례해 자라므로, 실제 오차는 대개 상한보다 훨씬 작다. 이 직관은 최근에 확률적 반올림 오차 해석으로 엄밀하게 정리되었다\cite{higham2019probabilistic}.

\subsection{나무 모양의 합}

\cref{thm:recursive-sum}의 증명에서 중요한 것은 항 $x_i$가 \textbf{몇 번의 덧셈을 거쳐} 최종 합에 이르는가이다. 더하는 순서를 이진 나무로 나타내면 각 잎이 항이고 각 안쪽 마디가 덧셈 하나인데, 같은 논리로 항 $x_i$의 상대 오차 인수는 잎에서 뿌리까지의 경로에 있는 덧셈의 수 $h_i$에 대해 $|\theta^{(i)}|\le\gamma_{h_i}$이다. 따라서
\[
|\hat s-s|\le\sum_{i=1}^n\gamma_{h_i}|x_i|\le\gamma_h\sum_{i=1}^n|x_i|,\qquad h=\max_ih_i
\]
이다. 재귀 합은 가장 한쪽으로 치우친 나무여서 $h=n-1$이다. 항을 둘씩 짝지어 더하고, 그 결과를 다시 둘씩 더하는 \term{쌍별 합}{pairwise summation}은 균형 잡힌 나무이므로 $h=\lceil\log_2n\rceil$이고, 오차의 상한이 $\gamma_{\lceil\log_2n\rceil}\sum|x_i|$로 줄어든다\cite{higham1993summation}. 백만 개의 합이라면 상한이 $10^6u$에서 $20u$로 줄어든다. 덧셈의 횟수는 그대로이고 순서만 바뀌었다.

\begin{instardust}[\code{csrc/base/pairwise.hpp}]
stardust의 모든 부동소수점 합은 이 파일의 나무 하나를 쓴다. 길이가 $512$ 이하인 구간은 잎에서 더하는데, $16$개의 독립된 누산기에 항을 번갈아 넣어 각자 차례로 더한 뒤 그 $16$개를 정해진 이진 나무로 모은다. 더 긴 구간은 $16$의 배수 경계에서 반으로 나누어 재귀적으로 더한다. 위의 방법으로 상한을 세어 보자. 잎 하나에서 항이 거치는 덧셈은 누산기의 사슬에서 $31$번 이하, 누산기를 모으는 나무에서 $4$번이고, 길이가 $16$의 배수가 아닌 잎에서는 마지막에 남은 $15$개 이하의 항을 차례로 더하는 몫이 붙으므로 모두 합쳐 $49$번을 넘지 않는다. 그 위로는 나누는 단계마다 덧셈이 하나씩이므로, 예컨대 $n=10^6$이면 $11$단계가 더해져 $h\le60$이고 오차의 상한은 $\gamma_{60}\sum|x_i|\approx6.7\times10^{-15}\sum|x_i|$이다. 잎에서 $16$개의 사슬을 따로 두는 이유는 정확도보다 속도에 있다. 서로 기다릴 필요가 없는 덧셈 $16$개가 NEON 레지스터 여덟 개를 채워 파이프라인을 쉬지 않게 하기 때문이며, 이 이야기는 13장에서 다시 한다. 나무의 모양이 길이에만 의존하므로 같은 길이의 합은 스레드 수와 상관없이 언제나 같은 비트를 낸다.
\end{instardust}

\subsection{내적}

$x^{\mathsf T}y=\sum_ix_iy_i$를 곱을 반올림한 뒤 차례로 더해 계산하면, 항 $x_iy_i$는 곱셈 한 번과 덧셈 많아야 $n-1$번을 거치므로
\[
\bigl|\fl(x^{\mathsf T}y)-x^{\mathsf T}y\bigr|\le\gamma_n\sum_{i=1}^n|x_iy_i|=\gamma_n\,|x|^{\mathsf T}|y|
\]
이다. 곱셈과 덧셈을 FMA로 합치면 인수가 하나 줄어들 뿐 꼴은 같고, 합하는 순서를 나무로 바꾸면 $\gamma_n$이 $\gamma_{\lceil\log_2n\rceil+1}$로 바뀐다. 행렬 곱의 각 성분은 내적이므로 같은 상한이 성분마다 성립하며, 이것은 19장 이후에 수치선형대수의 오차를 분석할 때 쓰는 기본 재료가 된다.

\begin{aside}[후진 오차 해석의 탄생]
컴퓨터가 막 등장한 1940년대 후반, 큰 행렬의 역행렬을 계산하는 것이 가능한지를 두고 비관론이 우세했다. 반올림 오차가 단계마다 쌓여 행렬의 크기에 대해 지수적으로 자란다는 어림이 있었기 때문이다. von Neumann과 Goldstine은 1947년의 긴 논문에서 양의 정부호 행렬의 역행렬에 대한 오차를 엄밀하게 묶어 이 비관론에 처음으로 답했고\cite{vonneumann1947}, 이듬해 Turing은 행렬의 조건수라는 말을 처음 쓰면서 오차를 문제의 성질과 알고리즘의 성질로 나누어 보는 관점을 내놓았다\cite{turing1948}. 이 생각을 체계적인 방법으로 완성한 사람은 영국 국립물리연구소에서 Turing과 함께 일했던 Wilkinson이다. 그는 계산 결과를 ``조금 바뀐 문제의 정확한 답''으로 해석하는 후진 오차 해석을 수많은 알고리즘에 적용해, Gauss 소거와 고유값 계산이 실제로는 아주 정확하다는 것을 보였다\cite{wilkinson1963rounding}.
\end{aside}

\section{소거}
\label{sec:cancellation}

\subsection{소거의 원리}

\cref{eq:cond-sub}에서 보았듯이 가까운 두 수의 뺄셈은 조건수가 크다. 8장의 Sterbenz 보조정리에 따르면 그 뺄셈 자체는 정확할 수도 있는데, 그래도 문제가 되는 이유는 뺄셈의 입력이 이미 오차를 품고 있기 때문이다. $\hat a=a(1+\Delta a)$, $\hat b=b(1+\Delta b)$이면 뺄셈이 정확하더라도
\[
\frac{|(\hat a-\hat b)-(a-b)|}{|a-b|}=\frac{|a\Delta a-b\Delta b|}{|a-b|}\le\max(|\Delta a|,|\Delta b|)\,\frac{|a|+|b|}{|a-b|}
\]
이고, 우변의 몫이 클수록 입력의 작은 상대오차가 결과에서 크게 자란다. 앞자리의 같은 숫자들이 지워지면서 뒷자리에 숨어 있던 오차가 앞으로 나온다. 이것을 \term{파국적 소거}{catastrophic cancellation}라 한다. 반대로 입력이 정확한 값이면, 예컨대 두 입력이 모두 사용자가 준 데이터라면, 뺄셈은 많아야 반올림 한 번의 오차만 만들므로 \term{온순한 소거}{benign cancellation}라 한다. 소거를 피하는 기술은 대부분 앞의 경우를 뒤의 경우로 바꾸는 것, 곧 오차가 생긴 중간값끼리 빼지 않도록 식을 고쳐 쓰는 것이다.

\subsection{식을 고쳐 쓰기}

예를 몇 가지 보자. $f(x)=(1-\cos x)/x^2$을 $x=1.2\times10^{-5}$에서 계산하면, $\cos x$의 binary64 값은 $0.999999999928$이고 공식대로 계산한 결과는 $0.49999973297\cdots$이다. 참값은 급수 $1/2-x^2/24+\cdots$에서 $0.499999999994$이므로 상대오차가 약 $5\times10^{-7}$, 곧 유효숫자 열 자리를 잃었다. $\cos x$는 상대오차 $u$ 이내로 정확하지만 그 오차 약 $10^{-16}$이 $1-\cos x\approx7.2\times10^{-11}$에 대한 상대오차로는 $10^{-6}$ 정도가 되기 때문이다. 반각 공식으로 $1-\cos x=2\sin^2(x/2)$라 고쳐 쓰면 $f(x)=2\bigl(\sin(x/2)/x\bigr)^2$이 되어 뺄셈이 사라지고, 이 식은 $0.499999999994$를 정확히 낸다.

$\log(1+x)$도 작은 $x$에서 같은 문제가 있다. $1+x$를 반올림하는 순간 $x$의 뒷자리가 사라지므로 $x=10^{-10}$에서 $\log(1\oplus x)$의 상대오차는 약 $8\times10^{-8}$이다. 그래서 수학 라이브러리는 $1+x$를 만들지 않고 $\log(1+x)$를 계산하는 \code{log1p}와, 같은 이유로 $e^x-1$을 직접 계산하는 \code{expm1}을 따로 제공한다. 8장에서 본 \code{expm1} 커널은 이 원리를 반올림 오차 없는 뺄셈과 함께 쓴 것이다.

이차방정식 $x^2+bx+c=0$의 근의 공식 $(-b\pm\sqrt{b^2-4c})/2$도 고전적인 예이다. $b=-10^8$, $c=1$이면 두 근은 약 $10^8$과 $10^{-8}$인데, 작은 근을 공식대로 $(-b-\sqrt{b^2-4c})/2$로 계산하면 거의 같은 두 수 $10^8$과 $\sqrt{10^{16}-4}$를 빼게 되어 $7.45\times10^{-9}$이라는, $25\%$나 틀린 값이 나온다. 큰 근은 같은 부호끼리 더하므로 소거가 없고, 두 근의 곱이 $c$라는 사실에서 작은 근을 $c/(\text{큰 근})$, 곧 $2c/(-b+\sqrt{b^2-4c})$로 얻으면 $1.0\times10^{-8}$이 정확히 나온다. 공식을 조금 바꾸었을 뿐인데 빼는 대상이 반올림된 중간값에서 사라졌다.

\subsection{분산}

통계에서 가장 유명한 소거의 예는 분산이다. 표본 $x_1,\dots,x_n$의 제곱 편차의 합 $S=\sum(x_i-\bar x)^2$은 대수적으로
\[
S=\sum_ix_i^2-\frac1n\Bigl(\sum_ix_i\Bigr)^2
\]
과 같아서 자료를 한 번만 훑고도 계산할 수 있으므로 오래전부터 교과서 공식으로 쓰였다. 그러나 이 식은 큰 두 수의 차이다. 데이터가 $10^8+\{4,7,13,16\}$이면 참값의 표본분산은 $S/(n-1)=90/3=30$인데, 교과서 공식은 $29.333\cdots$을 내고, 데이터를 $10^9+\{4,7,13,16\}$으로 옮기면 $0$을 낸다. 두 합 $\sum x_i^2\approx4\times10^{18}$이 반올림될 때의 오차가 이미 $S=90$보다 크기 때문이다.

이 현상은 조건수로 정확히 설명된다. Chan, Golub, LeVeque는 $S$의 데이터에 대한 조건수가
\[
\kappa=\sqrt{1+n\bar x^2/S}
\]
이고, 교과서 공식의 상대오차가 대략 $n\kappa^2u$, 평균을 먼저 구한 뒤 편차를 제곱해 더하는 \term{두 번 훑기}{two-pass} 알고리즘의 상대오차가 대략 $nu+(n\kappa u)^2$임을 보였다\cite{chan1983variance}. 위의 첫 예에서 $\kappa^2\approx4.4\times10^{14}$이므로 $n\kappa^2u\approx0.2$로 교과서 공식은 유효숫자를 거의 다 잃을 수 있고, 두 번 훑기는 $\kappa u\approx2.3\times10^{-9}$이 작으므로 정확하다. 교과서 공식의 오차가 $\kappa^2$에 비례하는 것은 데이터를 제곱하는 순간 조건수가 제곱되기 때문이며, 같은 현상을 22장에서 최소제곱의 정규방정식에서 다시 만난다. 자료를 한 번만 훑어야 할 때는 평균과 $S$를 새 자료가 올 때마다 갱신하는 Welford의 방법이 소거 없이 같은 일을 한다\cite{welford1962}.

\begin{instardust}[\code{docs/design/array.md} §6.8]
stardust의 \code{var}는 두 번 훑기이다. 첫 훑기에서 평균을 구하고, 둘째 훑기에서 평균과의 편차를 제곱해 더하되 임시 배열 없이 커널 하나로 접는다. 제곱을 더할 때는 \code{fma(d, d, acc)}로 곱과 덧셈을 반올림 한 번으로 묶는데, \code{pairwise.hpp}의 \code{leaf\_squares}가 그 잎이다. 통계 영역의 기술통계는 한 걸음 더 나아가, 결과가 오차를 키우는 합(예컨대 왜도를 위한 세제곱 편차의 합)만 골라 10장의 오차 없는 덧셈으로 오차까지 함께 더한다.
\end{instardust}

\subsection{정수로 소거를 없애기}

소거를 피하는 가장 확실한 방법은 빼는 두 양을 정확하게 만드는 것이다. 순위 검정의 통계량은 그 좋은 예이다. Kruskal--Wallis 검정의 통계량은 교과서에서 보통
\[
H=\frac{12}{N(N+1)}\sum_{i=1}^k\frac{R_i^2}{n_i}-3(N+1)
\]
로 적는다. 여기서 $R_i$는 $i$번째 표본의 순위합, $n_i$는 그 크기, $N=\sum n_i$이다. 두 항은 각각 $3(N+1)$ 정도의 크기이고 귀무가설 아래에서 그 차 $H$는 $k-1$ 정도이므로, $N$이 크면 큰 두 수의 차를 구하는 셈이 된다. 그런데 $\sum_iR_i=N(N+1)/2$이므로 편차 $D_i=2R_i-n_i(N+1)$을 쓰면
\[
\sum_i\frac{R_i^2}{n_i}-\frac{N(N+1)^2}4=\sum_i\frac{(R_i-n_i(N+1)/2)^2}{n_i}=\frac14\sum_i\frac{D_i^2}{n_i}
\]
이고, 따라서
\[
H=\frac{3}{N(N+1)}\sum_{i=1}^k\frac{D_i^2}{n_i}
\]
이다. $D_i$는 정수이므로 정확하게 계산되고, 이 식에는 빼는 항이 없다. 분산의 두 번 훑기가 평균을 먼저 빼듯이, 여기서도 평균에 해당하는 $n_i(N+1)/2$를 먼저 빼되 그 뺄셈을 정수에서 하는 것이다.

\begin{instardust}[\code{docs/design/htest.md} §3.4]
stardust의 Kruskal--Wallis와 Friedman 통계량은 정확한 정수 편차 $D_i$에서 위의 꼴로 계산한다(동점 보정이 든 식도 같은 방식으로 정리한다). 순위는 두 배 한 중간 순위로 들어 정수로 만들고(7장), $D_i^2$의 합은 double-double로 더한다. 설계 문서는 교과서 꼴을 쓰는 다른 구현이 큰 $N$에서 마지막 몇 자리를 잃는다고 적어 둔다.
\end{instardust}

\section{안정성}
\label{sec:stability}

알고리즘이 \term{수치적으로 안정}{numerically stable}하다는 말은 여러 뜻으로 쓰인다. 가장 강한 뜻은 후진 안정이고, 조금 약한 뜻은 결과가 ``조금 바뀐 입력에 대한 거의 정확한 답'', 곧 $\hat y+\Delta y=f(x+\Delta x)$에서 $|\Delta y|/|y|$와 $|\Delta x|/|x|$가 모두 작다는 \term{혼합 안정성}{mixed stability}이다. 어느 뜻이든 핵심은 알고리즘이 문제의 조건수 이상으로 오차를 키우지 않는다는 것이다. 반대로 \textbf{불안정한} 알고리즘은 조건이 좋은 문제에서도 큰 오차를 낸다.

6장의 Bessel 점화식은 불안정성의 교과서적인 예였다. $J_n(x)$의 값 자체는 조건이 나쁜 문제가 아닌데, 차수가 올라가는 방향으로 점화식을 돌리면 지배해가 반올림 오차를 증폭해 몇 걸음 만에 결과를 덮어 버린다. 같은 점화식을 거꾸로 돌리는 Miller의 알고리즘은 안정하다. 위에서 본 이차방정식과 분산의 예도 같은 구조로, 문제 자체의 조건은 좋거나 적당한데 특정한 계산 순서가 중간에 조건이 나쁜 단계를 만들어 넣은 것이다. 그래서 알고리즘의 안정성을 따질 때는 각 단계의 조건수를 보아야 한다. 어떤 단계의 조건수가 전체 문제의 조건수보다 훨씬 크면, 그 단계가 입력의 반올림 오차를 문제 자체가 요구하는 것보다 크게 키운다.

이 장의 도구로 알고리즘의 오차를 분석하는 과정은 대개 세 걸음이다. 표준 모형으로 각 연산의 반올림을 $(1+\delta)$로 적고, \cref{lem:gamma-n}으로 인수들을 $\theta_n$으로 묶어 계산 결과를 조금 바뀐 입력의 정확한 결과로 해석하며, 끝으로 문제의 조건수를 곱해 전진 오차를 얻는다. 상한은 대개 비관적이지만 오차가 어디서 생기고 무엇이 그것을 키우는지를 정확히 보여 주고, 그래서 알고리즘을 고치는 방향을 가리킨다.

\begin{instardust}[\code{docs/design/stats.md}]
확률분포의 역함수, 예컨대 정규분포의 분위수 $\mu+\sigma z$는 $\mu$와 $\sigma z$가 거의 같은 크기에 반대 부호이면 결과가 $0$ 근처로 소거된다. 이 소거는 알고리즘의 결함 탓이 아니라 문제의 조건이 나쁜 탓이어서 어떤 방법으로도 $z$의 반올림 오차가 $|\sigma z/(\mu+\sigma z)|$배로 커지는 것을 막을 수 없다. 그래서 stardust는 이런 역함수의 정확도를 ``$\max(1,\kappa)$ 단위로'' 약속한다. 설계 문서의 예에서 \code{Normal(1.5, 2).icdf(0.2267)}은 $4.824\times10^{-4}$이고 상대오차가 $3.3\times10^{-13}$이지만, 조건수가 약 $3120$이므로 계약의 단위로는 반올림 한 번 남짓의 오차이다.
\end{instardust}

\section{맺음}

반올림 오차 해석의 결론은 두 문장으로 줄일 수 있다. 알고리즘은 후진 오차로 평가하고, 결과의 정확도는 거기에 문제의 조건수를 곱해 어림한다. 그리고 오차가 커지는 곳은 거의 언제나 반올림된 중간값끼리의 소거이므로, 식을 고쳐 쓰거나 빼는 양을 정확하게 만들어 피한다. 다음 장에서는 한 걸음 더 나아가, 반올림 오차 자체를 정확하게 계산해 내서 정밀도를 두 배로 늘리는 방법을 다룬다.

이 장의 내용은 거의 모두 Higham의 책\cite{higham2002accuracy}에 더 자세하고 엄밀하게 정리되어 있으며, Trefethen과 Bau의 교과서\cite{trefethen1997nla}는 조건수와 안정성의 개념을 특히 명료하게 설명한다.

\chapterreferences
